
\documentclass{article}
\usepackage{colortbl}
\usepackage{makecell}
\usepackage{multirow}
\usepackage{supertabular}

\begin{document}

\newcounter{utterance}

\centering \large Interaction Transcript for game `ifeval', experiment `ifeval', episode 54 with score{-}Qwen3.5{-}9B{-}sp.
\vspace{24pt}

{ \footnotesize  \setcounter{utterance}{1}
\setlength{\tabcolsep}{0pt}
\begin{supertabular}{c@{$\;$}|p{.15\linewidth}@{}p{.15\linewidth}p{.15\linewidth}p{.15\linewidth}p{.15\linewidth}p{.15\linewidth}}
    \# & \multicolumn{2}{c}{Player} && \multicolumn{2}{c}{Game Master} \\
    \hline

    \theutterance \stepcounter{utterance}  
    & & & \multicolumn{4}{p{0.6\linewidth}}{
        \cellcolor[rgb]{0.9,0.9,0.9}{
            \makecell[{{p{\linewidth}}}]{
                \texttt{\tiny{[P1$\langle$GM]}}
                \texttt{Take the text below as a starting point, and make it a complete article: "You may have to meet with a helper to work out a parenting plan. The first would be to estimate how many time you have everyday for parenting, and is that enough...."} \\
\texttt{Avoid using the following keywords: sleep, cook, feed} \\
\texttt{Mention the keyword 'schedule' for more than 5 times.} \\
            }
        }
    }
     \\ \\

    \theutterance \stepcounter{utterance}  
    & \multicolumn{4}{p{0.6\linewidth}}{
        \cellcolor[rgb]{0.561,0.737,0.561}{
            \makecell[{{p{\linewidth}}}]{
                \texttt{\tiny{[P1$\rangle$GM]}}
                \texttt{\# Navigating the Path to a Parenting Plan with a Helper} \\
\\ 
\texttt{Creating a stable environment for your children often requires more than just love and intention; it frequently demands a structured approach, especially when you are seeking assistance from a helper. You may have to meet with a helper to work out a parenting plan that aligns with your family's unique needs. The first step in this collaborative process is to estimate how much time you have every day for parenting, and determine if that amount is sufficient to meet your children's developmental requirements.} \\
\\ 
\texttt{Once you have a clear understanding of your available hours, the next critical phase is drafting a detailed **schedule**. This **schedule** should outline specific blocks of time dedicated to various activities, ensuring that both you and the helper know exactly when to be present and when to step back. A well{-}organized **schedule** prevents confusion and reduces stress, allowing everyone involved to focus on the children rather than wondering about the next task.} \\
\\ 
\texttt{When establishing this **schedule**, it is essential to consider the rhythm of the day. You must decide how to divide responsibilities so that the **schedule** remains flexible enough to handle unexpected events while maintaining consistency. If the initial **schedule** proves too rigid or leaves gaps in coverage, it is important to revisit the plan with your helper. Adjustments might be necessary to ensure that the **schedule** supports the children's emotional and physical well{-}being without overwhelming the caregivers.} \\
\\ 
\texttt{Ultimately, the goal is to create a framework where the **schedule** serves as a reliable guide rather than a source of anxiety. By working closely with a helper to refine this **schedule**, you can build a supportive system that allows you to be the best parent possible, knowing that your time is being used effectively and efficiently.} \\
            }
        }
    }
    & & \\ \\

    \theutterance \stepcounter{utterance}  
    & & & \multicolumn{2}{p{0.3\linewidth}}{
        \cellcolor[rgb]{0.9,0.9,0.9}{
            \makecell[{{p{\linewidth}}}]{
                \texttt{\tiny{[GM$|$GM]}}
                {'keywords:forbidden_words': {'forbidden_words': ['sleep', 'cook', 'feed']}, 'keywords:frequency': {'relation': 'at least', 'keyword': 'schedule', 'frequency': 6}}
            }
        }
    }
    & & \\ \\

    \theutterance \stepcounter{utterance}  
    & & & \multicolumn{2}{p{0.3\linewidth}}{
        \cellcolor[rgb]{0.9,0.9,0.9}{
            \makecell[{{p{\linewidth}}}]{
                \texttt{\tiny{[GM$|$GM]}}
                \texttt{game\_result = WIN} \\
            }
        }
    }
    & & \\ \\

\end{supertabular}
}

\end{document}
